\selectlanguage{english}
\chapter{Introduction}

\section{Motivation}

The web of human knowledge is inherently a graph. Wikipedia alone encodes hundreds of
thousands of entities---people, places, organisations, concepts---and the hyperlinks
between their articles encode a rich relational structure that mirrors how ideas are
connected. Predicting which edges are missing from such a graph, a task known as
\textit{link prediction}, has applications ranging from knowledge-base completion and
recommender systems to citation network analysis and drug interaction discovery. The
challenge is compounded when nodes carry rich textual attributes: the graph topology
alone may be sparse or ambiguous, while the text alone may be too high-dimensional for
simple similarity measures. Jointly exploiting both sources of signal is therefore both
necessary and non-trivial.

Traditional link prediction methods fall into two broad families. Structural
heuristics---such as Common Neighbours, Adamic-Adar, and Preferential
Attachment---are computationally cheap and interpretable, but they are blind to textual
content and fail completely when two nodes share no common neighbours. Graph embedding
methods such as Node2Vec \cite{grover2016node2vec} learn low-dimensional representations
that capture neighbourhood topology, yet they still ignore node text and require
retraining whenever the graph changes. On the other side of the spectrum, large
pre-trained language models (PLMs) such as BERT \cite{devlin2019bert} and its
descendants can encode rich semantic relationships directly from text, and recent work
has demonstrated that fine-tuned sentence encoders can achieve near-perfect accuracy on
the DSAA 2023 Wikipedia link prediction benchmark \cite{phan2023link,tran2023text}.
However, running a transformer over every candidate pair is expensive at scale, and the
reported near-perfect scores raise questions about dataset characteristics---in
particular, whether trivially identifiable pairs such as self-loops or structurally
obvious links are inflating the aggregate metric.

This thesis is motivated by three observations that are underexplored in the existing
literature. First, the cost of full-transformer inference is prohibitive at the scale of
real-world graphs: the DSAA 2023 test set contains 238,000 node pairs, and encoding
every pair independently with a sentence transformer is orders of magnitude more
expensive than a heuristic lookup. Second, different subpopulations of node pairs
respond very differently to different feature families: structurally rich pairs are
trivially classified by topology, cold-start pairs with zero common neighbours require
semantic features, and genuinely ambiguous pairs may need deep language understanding.
Third, Kim et al.\ \cite{kim2024llms} showed that more than 78\% of new co-authorship
links involve at least one node with zero common neighbours with the other, highlighting
that cold-start handling is not an edge case but the dominant regime in dynamic graphs.

These observations motivate a \textit{tiered} approach: apply the cheapest classifier
first, escalate to more expensive classifiers only when the cheaper tier is uncertain,
and measure both accuracy and inference cost at each tier. Such a cascade aligns well
with the practical deployment scenario where a large fraction of pairs can be resolved
cheaply, and only a small residual requires heavy computation.

\section{Problem Statement}

Let $G = (V, E)$ be an undirected graph where each node $v \in V$ is associated with a
text document $d_v$ (in our setting, the Wikipedia article text). Given a set of node
pairs $\mathcal{Q} = \{(u, v)\}$, the \textit{link prediction} task is to assign a
binary label $y_{uv} \in \{0, 1\}$ to each pair, where $y_{uv} = 1$ indicates that an
edge $(u, v) \in E$ exists (or should exist) and $y_{uv} = 0$ indicates no edge. The
task is evaluated on held-out pairs whose true labels are withheld at inference time.

In the DSAA 2023 Wikipedia benchmark \cite{phan2023link}, the graph $G$ has 838,000
nodes corresponding to Wikipedia articles, of which 668,000 appear in the structural
(adjacency) component. Training labels are provided for 948,232 node pairs (46\%
positive, 54\% negative), and the evaluation metric is the Macro F1-score computed over
the binary positive and negative classes. The test set contains 238,000 pairs. Node text
is available as a \texttt{nodes.tsv} file containing Wikipedia titles, descriptions, and
optionally infobox content.

The specific problem addressed by this thesis is to build a link prediction system for
this benchmark that (a) achieves competitive Macro F1-score, (b) explicitly manages
inference cost by routing easy pairs away from expensive models, (c) is transparent about
which feature family is responsible for each prediction, and (d) handles cold-start pairs
where one or both nodes have no structural neighbours in $G$. The framework must operate
at inference time without access to the true labels of the test pairs, using only the
training graph, the node text, and the training pair labels.

A secondary objective is to characterise the dataset itself: to identify which
subpopulations of pairs are trivially predictable (e.g., self-loops, pairs with many
common neighbours) and to quantify what fraction of the test set genuinely requires
semantic reasoning. This characterisation informs the threshold tuning for the cascade
and contextualises the near-perfect scores reported by prior work.

\section{Contributions}

The main contributions of this thesis are:

\begin{enumerate}

    \item \textbf{CascadeLP, a three-tier confidence-based link prediction framework.}
          We design and implement a cascade in which Tier~0 resolves self-loops
          deterministically, Tier~1 applies a Logistic Regression classifier over four
          structural heuristic features (Common Neighbours, Jaccard Coefficient,
          Adamic-Adar Index, and Preferential Attachment), Tier~2 applies a Random
          Forest classifier over a 72-dimensional part-of-speech frequency feature
          vector derived from the node text, and Tier~3 applies a Logistic Regression
          classifier over the cosine similarity computed by a Sentence-Transformer
          encoder. Pairs escalate to the next tier only when the current tier's maximum
          predicted probability falls below a tuned confidence threshold.

    \item \textbf{A systematic feature engineering study across structural, lexical,
          and semantic feature families.} We evaluate each feature family in isolation
          and in the cascade setting, providing quantitative evidence for which feature
          family dominates which subpopulation of pairs. In particular, we demonstrate
          that structural features suffice for the majority of training pairs that have
          at least one common neighbour, while POS-tag features provide meaningful
          signal for cold-start pairs whose structural features collapse to a single
          degenerate value.

    \item \textbf{A dataset characterisation of the DSAA 2023 Wikipedia benchmark.}
          We identify 10,221 self-loop pairs (all positive) that are trivially
          predictable before any model is invoked, analyse the distribution of Common
          Neighbour counts across positive and negative training pairs, and quantify
          the fraction of pairs that fall into the cold-start regime. This
          characterisation explains why near-perfect aggregate scores are achievable on
          this benchmark and provides context for interpreting published results.

    \item \textbf{A confidence-threshold routing mechanism with per-tier calibration.}
          We design a routing algorithm that consults each tier's calibrated
          \texttt{predict\_proba} output and escalates to the next tier when confidence
          is below a threshold. We describe the threshold selection procedure and
          analyse the sensitivity of the overall F1-score and transformer call rate to
          threshold choice.

    \item \textbf{An end-to-end implementation on the DSAA 2023 benchmark with
          reproducible experiments.} The full pipeline---preprocessing, feature
          extraction, model training, cascade routing, and evaluation---is implemented
          in Python and made publicly available. We report Macro F1-score, per-class
          precision and recall, inference throughput, and the fraction of pairs
          resolved at each tier, providing a concrete efficiency-accuracy trade-off
          analysis that prior work has not attempted.

\end{enumerate}

\section{Thesis Structure}

The remainder of this thesis is organised as follows.

Chapter~\ref{chap:background} surveys the relevant literature across the three main
paradigms for link prediction in text-attributed graphs: structural heuristics and graph
embeddings, pre-trained language models and sentence encoders, and hybrid architectures
that combine graph topology with text. It also describes the DSAA 2023 Wikipedia
benchmark in detail and identifies the gaps in existing work that motivate CascadeLP.
The chapter draws on foundational work including the Adamic-Adar index
\cite{adamic2003friends}, the Barab{\'a}si-Albert preferential attachment model
\cite{barabasi1999emergence}, Node2Vec \cite{grover2016node2vec}, BERT
\cite{devlin2019bert}, RoBERTa \cite{liu2019roberta}, Sentence-BERT
\cite{reimers2019sentence}, LinkBERT \cite{yasunaga2022linkbert}, GraphFormers
\cite{yang2021graphformers}, and Patton \cite{jin2023patton}.

Chapter~\ref{chap:methodology} presents the CascadeLP framework in full technical
detail. It covers the data preprocessing pipeline (text cleaning, infobox and
description extraction), the feature engineering for each tier (heuristic indices,
POS-tag frequency vectors, Sentence-Transformer cosine similarity), the training of
each tier's classifier, the confidence-based routing algorithm, and the cold-start
handling strategy.

Chapter~4 describes the experimental setup: the train/validation split strategy,
hyperparameter tuning, evaluation protocol, and baseline models used for comparison.
It reports the main quantitative results on the DSAA 2023 test set, including per-tier
accuracy, aggregate Macro F1-score, and inference cost.

Chapter~5 presents ablation studies that isolate the contribution of each tier and
each feature family, a cold-start analysis that evaluates performance on the subset of
pairs with zero common neighbours, and an efficiency benchmark that measures
throughput and the fraction of pairs escalated to Tier~3.

Chapter~6 summarises the findings, discusses limitations of the current framework,
and outlines directions for future work including dynamic threshold adaptation, graph
neural network integration, and evaluation on additional text-attributed graph
benchmarks.
